% GENERATED FILE. Edit canonical lessons, then run npm run generate:books.
% canonical-source-hash: fnv1a64:17e2055b9e1946a2
% canonical-lessons: HI-C464-ubna, HI-C464-dilcaspi, HI-C464-utsuk, HI-C464-jijnasa, HI-C464-rocak

\chapter{Bored or Curious}
\label{ch:hi-bored-or-curious}
\hlchaptermodality{464}

\begin{chapteropening}
\textbf{By the end of this chapter:} \emph{I can say what bores me and what interests me.}

Complete the last lesson of chapter 464: I can say what bores me and what interests me.
\end{chapteropening}

\section[\texorpdfstring{\hi{ऊबना} (ūbnā)}{ऊबना (ūbnā)}]{\hi{ऊबना} (ūbnā) — to get bored}

\label{lesson:HI-C464-ubna}



\begin{quote}
Before the new one: say the Hindi for to eat one's fill, then the Hindi for consolation.
\end{quote}

\subsection*{You'll want to know: \hi{ऊबना}}
\textbf{\hi{ऊबना}} — \emph{ūbnā} — \textquotedblleft{}to get bored\textquotedblright{}.

\begin{quote}
\hi{रोज़} \hi{दाल} \hi{खाकर} \hi{अरुण} \hi{ऊब} \hi{गया।}
\end{quote}

\begin{quote}
\emph{roz dāl khākar Aruṇ ūb gayā.}
\end{quote}

\begin{quote}
Arun got bored of eating dal every day.
\end{quote}

Past (he …): \hi{वह} \hi{ऊब} \hi{गया}, \emph{vah ūb gayā}.

\begin{cousinweb}
\textbf{\hi{ऊबना}} (\emph{ūbnā}) is inherited: CDIAL lists Hindi \emph{ūbnā}, with the sense \textquotedblleft{}to be agitated, be bewildered\textquotedblright{}, under Sanskrit \emph{údvijatē}, \textquotedblleft{}is agitated\textquotedblright{} (CDIAL 2087). The modern sense, the restless weariness of boredom, is a later turn that CDIAL does not give.
\end{cousinweb}

\begin{grammarlens}[title={What you've built}]
One more word for describing what you lived and felt.
\end{grammarlens}

\subsection*{Your turn}
\begin{itemize}
\raggedright
  \item \emph{Say it:} \emph{ūbnā}
  \item \emph{Say it:} \emph{ūbnā}, once more
  \item \emph{Say it:} \emph{roz dāl khākar Aruṇ ūb gayā.} — Arun got bored of eating dal every day.
  \item \emph{Recall:} say the Hindi for to eat one's fill, then the Hindi for consolation, then say \emph{ūbnā} again
\end{itemize}

\begin{tcolorbox}[breakable,colback=teal!4,colframe=teal!35!black,title={Before you move on}]
What does \hi{ऊबना} mean? (\textquotedblleft{}To get bored\textquotedblright{}.) Say it once more.
\end{tcolorbox}
\section[\texorpdfstring{\hi{दिलचस्पी} (dilcaspī)}{दिलचस्पी (dilcaspī)}]{\hi{दिलचस्पी} (dilcaspī) — interest; a keen interest}

\label{lesson:HI-C464-dilcaspi}



\begin{quote}
Before the new one: say the Hindi for consolation, then the Hindi for to get bored.
\end{quote}

\subsection*{You'll want to know: \hi{दिलचस्पी}}
\textbf{\hi{दिलचस्पी}} — \emph{dilcaspī} — \textquotedblleft{}interest; a keen interest\textquotedblright{}.

\begin{quote}
\hi{मीरा} \hi{को} \hi{इतिहास} \hi{में} \hi{बहुत} \hi{दिलचस्पी} \hi{है।}
\end{quote}

\begin{quote}
\emph{Mīrā ko itihās meṁ bahut dilcaspī hai.}
\end{quote}

\begin{quote}
Mira is very interested in history.
\end{quote}

Feminine.

\begin{cousinweb}
\textbf{\hi{दिलचस्पी}} (\emph{dilcaspī}) is usually traced to Persian \emph{dilchaspī}, \textquotedblleft{}interest\textquotedblright{}, built on \emph{dil}, \textquotedblleft{}heart\textquotedblright{}, and a word for \textquotedblleft{}sticking\textquotedblright{}: what the heart sticks to.
\end{cousinweb}

\begin{grammarlens}[title={What you've built}]
One more word for describing what you lived and felt.
\end{grammarlens}

\subsection*{Your turn}
\begin{itemize}
\raggedright
  \item \emph{Say it:} \emph{dilcaspī}
  \item \emph{Say it:} \emph{dilcaspī}, once more
  \item \emph{Say it:} \emph{Mīrā ko itihās meṁ bahut dilcaspī hai.} — Mira is very interested in history.
  \item \emph{Recall:} say the Hindi for consolation, then the Hindi for to get bored, then say \emph{dilcaspī} again
\end{itemize}

\begin{tcolorbox}[breakable,colback=teal!4,colframe=teal!35!black,title={Before you move on}]
What does \hi{दिलचस्पी} mean? (\textquotedblleft{}Interest; a keen interest\textquotedblright{}.) Say it once more.
\end{tcolorbox}
\section[\texorpdfstring{\hi{उत्सुक} (utsuk)}{उत्सुक (utsuk)}]{\hi{उत्सुक} (utsuk) — eager; keen}

\label{lesson:HI-C464-utsuk}



\begin{quote}
Before the new one: say the Hindi for to get bored, then the Hindi for interest.
\end{quote}

\subsection*{You'll want to know: \hi{उत्सुक}}
\textbf{\hi{उत्सुक}} — \emph{utsuk} — \textquotedblleft{}eager; keen\textquotedblright{}.

\begin{quote}
\hi{बच्चे} \hi{नया} \hi{खिलौना} \hi{देखने} \hi{के} \hi{लिए} \hi{उत्सुक} \hi{थे।}
\end{quote}

\begin{quote}
\emph{bacce nayā khilaunā dekhne ke liye utsuk the.}
\end{quote}

\begin{quote}
The children were eager to see the new toy.
\end{quote}

\begin{cousinweb}
\textbf{\hi{उत्सुक}} (\emph{utsuk}) is a tatsama, borrowed from Sanskrit \emph{utsuka}, \textquotedblleft{}restless, uneasy\textquotedblright{} and so \textquotedblleft{}eager\textquotedblright{} (Monier-Williams).
\end{cousinweb}

\begin{grammarlens}[title={What you've built}]
One more word for describing what you lived and felt.
\end{grammarlens}

\subsection*{Your turn}
\begin{itemize}
\raggedright
  \item \emph{Say it:} \emph{utsuk}
  \item \emph{Say it:} \emph{utsuk}, once more
  \item \emph{Say it:} \emph{bacce nayā khilaunā dekhne ke liye utsuk the.} — The children were eager to see the new toy.
  \item \emph{Recall:} say the Hindi for to get bored, then the Hindi for interest, then say \emph{utsuk} again
\end{itemize}

\begin{tcolorbox}[breakable,colback=teal!4,colframe=teal!35!black,title={Before you move on}]
What does \hi{उत्सुक} mean? (\textquotedblleft{}Eager; keen\textquotedblright{}.) Say it once more.
\end{tcolorbox}
\section[\texorpdfstring{\hi{जिज्ञासा} (jijñāsā)}{जिज्ञासा (jijñāsā)}]{\hi{जिज्ञासा} (jijñāsā) — curiosity; the wish to know}

\label{lesson:HI-C464-jijnasa}



\begin{quote}
Before the new one: say the Hindi for interest, then the Hindi for eager.
\end{quote}

\subsection*{You'll want to know: \hi{जिज्ञासा}}
\textbf{\hi{जिज्ञासा}} — \emph{jijñāsā} — \textquotedblleft{}curiosity; the wish to know\textquotedblright{}.

\begin{quote}
\hi{छोटे} \hi{बच्चों} \hi{में} \hi{बहुत} \hi{जिज्ञासा} \hi{होती} \hi{है।}
\end{quote}

\begin{quote}
\emph{choṭe baccoṁ meṁ bahut jijñāsā hotī hai.}
\end{quote}

\begin{quote}
Small children are very curious.
\end{quote}

Feminine.

\begin{cousinweb}
\textbf{\hi{जिज्ञासा}} (\emph{jijñāsā}) is a tatsama, borrowed from Sanskrit \emph{jijñāsā}, \textquotedblleft{}the desire to know\textquotedblright{}, the desiderative of \emph{jñā}, \textquotedblleft{}to know\textquotedblright{} (Monier-Williams).
\end{cousinweb}

\begin{grammarlens}[title={What you've built}]
One more word for describing what you lived and felt.
\end{grammarlens}

\subsection*{Your turn}
\begin{itemize}
\raggedright
  \item \emph{Say it:} \emph{jijñāsā}
  \item \emph{Say it:} \emph{jijñāsā}, once more
  \item \emph{Say it:} \emph{choṭe baccoṁ meṁ bahut jijñāsā hotī hai.} — Small children are very curious.
  \item \emph{Recall:} say the Hindi for interest, then the Hindi for eager, then say \emph{jijñāsā} again
\end{itemize}

\begin{tcolorbox}[breakable,colback=teal!4,colframe=teal!35!black,title={Before you move on}]
What does \hi{जिज्ञासा} mean? (\textquotedblleft{}Curiosity; the wish to know\textquotedblright{}.) Say it once more.
\end{tcolorbox}
\section[\texorpdfstring{\hi{रोचक} (rocak)}{रोचक (rocak)}]{\hi{रोचक} (rocak) — interesting; engaging}

\label{lesson:HI-C464-rocak}



\begin{quote}
Before the new one: say the Hindi for eager, then the Hindi for curiosity.
\end{quote}

\subsection*{You'll want to know: \hi{रोचक}}
\textbf{\hi{रोचक}} — \emph{rocak} — \textquotedblleft{}interesting; engaging\textquotedblright{}.

\begin{quote}
\hi{दादा} \hi{की} \hi{कहानियाँ} \hi{बहुत} \hi{रोचक} \hi{होती} \hi{हैं।}
\end{quote}

\begin{quote}
\emph{dādā kī kahāniyā̃ bahut rocak hotī haiṁ.}
\end{quote}

\begin{quote}
Grandfather's stories are always very interesting.
\end{quote}

\begin{cousinweb}
\textbf{\hi{रोचक}} (\emph{rocak}) is a tatsama, borrowed from Sanskrit \emph{rocaka}, \textquotedblleft{}giving an appetite\textquotedblright{} and \textquotedblleft{}pleasing, agreeable\textquotedblright{} (Monier-Williams): what gives the mind an appetite.
\end{cousinweb}

\begin{grammarlens}[title={What you've built}]
One more word for describing what you lived and felt.
\end{grammarlens}

\subsection*{Your turn}
\begin{itemize}
\raggedright
  \item \emph{Say it:} \emph{rocak}
  \item \emph{Say it:} \emph{rocak}, once more
  \item \emph{Say it:} \emph{dādā kī kahāniyā̃ bahut rocak hotī haiṁ.} — Grandfather's stories are always very interesting.
  \item \emph{Recall:} say the Hindi for eager, then the Hindi for curiosity, then say \emph{rocak} again
\end{itemize}

\begin{tcolorbox}[breakable,colback=teal!4,colframe=teal!35!black,title={Before you move on}]
What does \hi{रोचक} mean? (\textquotedblleft{}Interesting; engaging\textquotedblright{}.) Say it once more.
\end{tcolorbox}
